\chapter{Data, प्रयोग एवं AI उपकरण}
\label{ch:data}

\lakshyaheading
\begin{itemize}[leftmargin=1.4em]
  \item अच्छे data (good data) के गुण - शुद्धता, विविधता एवं पर्याप्त मात्रा।
  \item डेटा-विज्ञान (data science) के पाँच चरण - संग्रह से निष्कर्ष तक।
  \item data में \term{पक्षपात}{bias} को पहचानना एवं उससे बचना।
  \item विद्यार्थियों के लिए मुफ़्त एवं open-source AI उपकरण।
  \item \emph{Google Colab} में अपना पहला Python कोड स्वयं चलाना।
\end{itemize}

\section{``Garbage in, garbage out''}

कंप्यूटर विज्ञान का एक प्रसिद्ध कथन है - \emph{``Garbage in, garbage out''}, अर्थात्
यदि भीतर डाला गया data ख़राब है तो परिणाम भी ख़राब ही होगा। कोई भी AI उतनी ही अच्छी
होती है जितना उसका डेटा। इसीलिए वैज्ञानिक प्रयोग में data का \emph{सावधानीपूर्वक संग्रह}
सबसे महत्वपूर्ण चरण माना जाता है।

अच्छे data के तीन प्रमुख गुण हैं:
\begin{itemize}[leftmargin=1.4em]
  \item \term{शुद्धता}{Accuracy} - मान सही ढंग से एवं सही उपकरण से मापे गए हों।
  \item \term{विविधता}{Diversity} - data सभी प्रकार की स्थितियों को दर्शाता हो।
  \item \term{पर्याप्त मात्रा}{Sufficient quantity} - इतने उदाहरण हों कि pattern
        स्पष्ट दिखे।
\end{itemize}

\section{डेटा-विज्ञान के पाँच चरण}

किसी भी जाँच में - चाहे भौतिकी की हो, रसायन की या जीव विज्ञान की - data के साथ कार्य
प्रायः इन्हीं पाँच चरणों में होता है:

\begin{enumerate}[leftmargin=1.8em]
  \item \term{संग्रह}{Collection} - सेंसर, सर्वेक्षण या प्रयोग से data एकत्र करना।
  \item \term{सफ़ाई}{Cleaning} - त्रुटियाँ, दोहराव एवं रिक्त (missing) मान हटाना।
  \item \term{विश्लेषण}{Analysis} - ग्राफ़ एवं सांख्यिकी (statistics) से pattern देखना।
  \item \term{मॉडलन}{Modelling} - pattern से नियम/पूर्वानुमान निकालना (जैसे regression)।
  \item \term{निष्कर्ष}{Conclusion} - परिणाम की वैज्ञानिक व्याख्या एवं सत्यापन।
\end{enumerate}

ये पाँच चरण प्रायः एक \emph{चक्र (cycle)} की तरह दोहराए जाते हैं - निष्कर्ष से नए
प्रश्न उठते हैं और फिर नया संग्रह आरम्भ होता है (चित्र~\ref{fig:ds-cycle})।

\begin{figure}[h]
\centering
\begin{tikzpicture}[font=\sffamily\small,
    dstep/.style={draw=aiblue, thick, fill=ailite, circle, align=center,
      minimum size=1.9cm, inner sep=1pt},
    arr/.style={-{Stealth[length=2.6mm]}, thick, aigreen}]
  \def\R{3cm}
  \node[dstep] (s1) at (90:\R)  {संग्रह\\\footnotesize Collection};
  \node[dstep] (s2) at (18:\R)  {सफ़ाई\\\footnotesize Cleaning};
  \node[dstep] (s3) at (-54:\R) {विश्लेषण\\\footnotesize Analysis};
  \node[dstep] (s4) at (-126:\R){मॉडलन\\\footnotesize Modelling};
  \node[dstep] (s5) at (162:\R) {निष्कर्ष\\\footnotesize Conclusion};
  \draw[arr] (s1) to[bend left=12] (s2);
  \draw[arr] (s2) to[bend left=12] (s3);
  \draw[arr] (s3) to[bend left=12] (s4);
  \draw[arr] (s4) to[bend left=12] (s5);
  \draw[arr] (s5) to[bend left=12] (s1);
\end{tikzpicture}
\caption{डेटा-विज्ञान का चक्र: संग्रह से निष्कर्ष तक, और फिर नए प्रश्नों की ओर।}
\label{fig:ds-cycle}
\end{figure}

\begin{jante}
data वैज्ञानिक (data scientists) प्रायः अपना 60-80\% समय केवल \emph{चरण~2 (cleaning)}
में लगाते हैं। असली कौशल अक्सर ``साफ़ data'' बनाने में ही होता है, न कि केवल model चलाने में।
\end{jante}

\section{data में पक्षपात (Bias)}

यदि एकत्र किया गया data एकपक्षीय हो, तो AI के परिणाम भी \term{पक्षपाती}{biased} होंगे -
भले ही कोड बिल्कुल सही हो। उदाहरण के लिए, यदि किसी पत्ती-पहचान model को केवल हरी पत्तियों
के चित्र दिए जाएँ, तो वह सूखी या पीली पत्ती पहचानने में असफल रहेगा।

इसीलिए वैज्ञानिक कार्य में यह सदैव पूछा जाता है - ``क्या मेरा data सभी स्थितियों का
प्रतिनिधित्व करता है?'' यह प्रश्न आगे अध्याय~6 (नैतिकता) से भी गहराई से जुड़ा है।

\begin{sochiye}
मान लीजिए किसी कक्षा की ``औसत लम्बाई'' केवल खेल-दल (sports team) के विद्यार्थियों से
मापी जाए। यह माप पूरी कक्षा के लिए सही क्यों नहीं होगी? यह डेटा-पक्षपात का उदाहरण कैसे है?
\end{sochiye}

\section{विद्यार्थियों के लिए उपयोगी मुफ़्त उपकरण}

इस पुस्तक की सभी गतिविधियाँ निःशुल्क (free) एवं प्रायः open-source उपकरणों से की जा
सकती हैं:

\begin{table}[h]
\centering
\begin{tabular}{@{}p{3.6cm} p{9.4cm}@{}}
\toprule
\rowcolor{cvBlue}\hdc{उपकरण (tool)} & \hdc{उपयोग} \\
\midrule
\href{https://teachablemachine.withgoogle.com/}{Teachable Machine} & बिना कोड के image/sound classification \\
\href{https://colab.research.google.com/}{Google Colab}      & ब्राउज़र में मुफ़्त Python चलाना (कुछ इंस्टॉल नहीं करना पड़ता) \\
\href{https://phyphox.org/}{Phyphox}           & मोबाइल सेंसर से भौतिक data एकत्र करना \\
\href{https://orangedatamining.com/}{Orange Data Mining} & दृश्य (visual) रूप से ML समझना, बिना अधिक कोड \\
\href{https://molview.org/}{MolView} / \href{https://pubchem.ncbi.nlm.nih.gov/}{PubChem}  & अणुओं की 3D संरचना देखना \\
\bottomrule
\end{tabular}
\caption{इस पुस्तक में प्रयुक्त मुफ़्त एवं प्रायः open-source उपकरण (नाम पर क्लिक करने योग्य)}
\end{table}

\noindent\footnotesize\textcolor{aigrey}{सूचना: ऊपर तालिका में प्रत्येक उपकरण का नाम
क्लिक करने योग्य लिंक है; डिजिटल (PDF) संस्करण में सीधे वेबसाइट खुल जाएगी।}\normalsize

\section{अपना पहला Python कोड (Google Colab में)}

पिछले अध्यायों में हमने कई बार Python कोड देखा। अब उसे स्वयं चलाने का समय है। \emph{Google
Colab} एक मुफ़्त सेवा है जो ब्राउज़र में ही Python चला देती है - कंप्यूटर पर कुछ भी
इंस्टॉल करने की आवश्यकता नहीं।

\begin{gatividhi}[Colab में पहला कदम - data का ग्राफ़ एवं regression]
\textbf{लक्ष्य:} स्वयं Python चलाकर data पर सर्वोत्तम रेखा खींचना।\\[4pt]
\textbf{सामग्री:} कंप्यूटर/मोबाइल, इंटरनेट, एक Google खाता।\\[4pt]
\textbf{चरण:}
\begin{enumerate}[leftmargin=1.6em]
  \item ब्राउज़र में \texttt{colab.research.google.com} खोलें और ``New Notebook'' चुनें।
  \item नीचे दिया कोड एक cell में लिखें (या copy करें)।
  \item cell के बाएँ ``play'' बटन (या Shift+Enter) दबाकर कोड चलाएँ।
  \item ग्राफ़ एवं छपा हुआ (printed) मान देखें।
\end{enumerate}
\end{gatividhi}

\begin{center}
\fbox{\qrcode[height=2.2cm]{https://colab.research.google.com/}}\\[3pt]
{\footnotesize\sffamily\color{aigrey} Google Colab खोलने हेतु इस QR को स्कैन करें\\
(\nolinkurl{colab.research.google.com})}
\end{center}

\begin{lstlisting}
import numpy as np
import matplotlib.pyplot as plt

# experiment data (e.g. a spring: force vs extension)
x = np.array([0.02, 0.04, 0.06, 0.08, 0.10])
y = np.array([1.0, 2.1, 2.9, 4.2, 5.0])

# modelling: best-fit straight line y = m*x + c
m, c = np.polyfit(x, y, 1)
print("slope m =", round(m, 2), " intercept c =", round(c, 2))

# analysis: graph of the data and the fitted line
plt.scatter(x, y, label="data")          # original data points
plt.plot(x, m*x + c, color="red", label="best-fit line")
plt.xlabel("x"); plt.ylabel("y"); plt.legend()
plt.show()
\end{lstlisting}
\colabnb{https://colab.research.google.com/github/kkreboot/ai-in-science-hindi/blob/main/notebooks/ch5-regression-plot.ipynb}

\noindent यह वही \emph{least-squares regression} है जिसे आपने अध्याय~2 में हाथ से किया
था - अब कंप्यूटर इसे कुछ ही क्षणों में, अधिक सटीकता से कर देता है।

\begin{gatividhi}[कक्षा का data - सहसम्बन्ध (correlation)]
\textbf{लक्ष्य:} वास्तविक data एकत्र कर उसमें सम्बन्ध खोजना।\\[4pt]
\textbf{चरण:}
\begin{enumerate}[leftmargin=1.6em]
  \item पूरी कक्षा की लम्बाई (height) एवं भुजा-विस्तार (arm-span) मापें।
  \item दोनों का scatter-plot बनाएँ (कागज़ पर या Colab में)।
  \item देखें - क्या लम्बाई बढ़ने पर भुजा-विस्तार भी बढ़ता है (correlation)?
\end{enumerate}
\textbf{नोट:} यहाँ ``height'' एवं ``arm-span'' को ML की भाषा में \emph{feature} कहते हैं।
\end{gatividhi}

% ---------------- परियोजना ----------------
\begin{pariyojana}[अपना पहला डेटासेट - पूरा डेटा-विज्ञान चक्र]
\textbf{उद्देश्य:} किसी एक प्रश्न पर स्वयं data एकत्र कर पाँचों चरणों का पालन करना।\\[4pt]
\textbf{कार्य (समूह में, 3-4 विद्यार्थी):}
\begin{enumerate}[leftmargin=1.6em]
  \item एक प्रश्न चुनें - जैसे ``दिन के विभिन्न समय पर कक्षा का तापमान कैसे बदलता है?''
        या ``पौधे की वृद्धि प्रतिदिन कितनी होती है?''
  \item \textbf{संग्रह:} कम से कम एक सप्ताह नियमित रूप से मान मापें एवं तालिका में लिखें।
  \item \textbf{सफ़ाई:} ग़लत या छूटे हुए मानों को जाँचें एवं ठीक करें।
  \item \textbf{विश्लेषण:} data का ग्राफ़ बनाएँ (कागज़ या Colab)।
  \item \textbf{मॉडलन:} यदि सम्भव हो तो एक प्रवृत्ति-रेखा (trend line) खींचें एवं किसी
        अगले मान का पूर्वानुमान करें।
  \item \textbf{निष्कर्ष:} परिणाम को 5 स्लाइड या एक पोस्टर के रूप में प्रस्तुत करें,
        और data की किन्हीं सीमाओं (limitations) का उल्लेख करें।
\end{enumerate}
\textbf{मूल्यांकन:} पाँचों चरणों के पालन, data की गुणवत्ता एवं निष्कर्ष की स्पष्टता
के आधार पर।
\end{pariyojana}

% ---------------- सारांश ----------------
\begin{bharatai}
\textbf{UPI} जैसी डिजिटल भुगतान प्रणालियाँ प्रतिदिन करोड़ों लेन-देन पर AI से धोखाधड़ी
(fraud) तुरंत पकड़ती हैं - यह वास्तविक-समय (real-time) डेटा-विश्लेषण का बेहतरीन उदाहरण है।
\end{bharatai}

\section*{सारांश (Summary)}
\addcontentsline{toc}{section}{सारांश}
\begin{itemize}[leftmargin=1.4em]
  \item अच्छा data (शुद्ध, विविध एवं पर्याप्त) ही अच्छी AI की नींव है -
        ``garbage in, garbage out''।
  \item डेटा-विज्ञान के पाँच चरण हैं: संग्रह, सफ़ाई, विश्लेषण, मॉडलन एवं निष्कर्ष।
  \item एकपक्षीय data से पक्षपाती (biased) परिणाम मिलते हैं, चाहे कोड सही हो।
  \item Teachable Machine, Google Colab, Phyphox एवं Orange जैसे मुफ़्त उपकरणों से
        विद्यार्थी स्वयं AI का अनुभव कर सकते हैं।
  \item Google Colab में कुछ ही पंक्तियों के Python से data का ग्राफ़ एवं regression
        किया जा सकता है।
\end{itemize}

% ---------------- अभ्यास ----------------
\section*{अभ्यास (Exercises)}
\addcontentsline{toc}{section}{अभ्यास}
\textbf{अ. रिक्त स्थान भरिए:}
\begin{enumerate}[leftmargin=1.6em]
  \item ``Garbage in, \rule{2.8cm}{0.4pt} out'' - ख़राब data से ख़राब परिणाम।
  \item डेटा-विज्ञान का वह चरण जिसमें त्रुटियाँ हटाई जाती हैं, \rule{2.8cm}{0.4pt} कहलाता है।
  \item ब्राउज़र में मुफ़्त Python चलाने का उपकरण \rule{2.8cm}{0.4pt} है।
\end{enumerate}

\textbf{ब. लघु उत्तरीय:}
\begin{enumerate}[leftmargin=1.6em]
  \item अच्छे data के तीन गुण लिखिए।
  \item डेटा-विज्ञान के पाँच चरण क्रम से बताइए।
  \item डेटा-पक्षपात (bias) क्या है? एक उदाहरण दीजिए।
\end{enumerate}

\textbf{स. क्रियात्मक / दीर्घ उत्तरीय:}
\begin{enumerate}[leftmargin=1.6em]
  \item Google Colab में ऊपर दिया कोड चलाइए तथा प्राप्त ढलान $m$ एवं $c$ के मान लिखिए।
  \item कक्षा के height बनाम arm-span data का scatter-plot बनाइए और सम्बन्ध का वर्णन कीजिए।
  \item किसी एक प्रश्न पर सप्ताह-भर data एकत्र कर पाँचों चरणों की एक संक्षिप्त रिपोर्ट
        बनाइए।
\end{enumerate}

\begin{mcq}
  \item ``Garbage in, garbage out'' का अर्थ है:
    \begin{choices}\item तेज़ कंप्यूटर \item ख़राब data से ख़राब परिणाम
      \item मुफ़्त सॉफ़्टवेयर \item बड़ा model\end{choices}
  \item ब्राउज़र में निःशुल्क Python चलाने की सेवा है:
    \begin{choices}\item Google Colab \item MS Word \item Photoshop \item VLC\end{choices}
  \item मोबाइल सेंसर से भौतिक data एकत्र करने वाला ऐप है:
    \begin{choices}\item Phyphox \item Spotify \item WhatsApp \item Zoom\end{choices}
  \item data में किसी वर्ग का अनुचित कम/अधिक प्रतिनिधित्व कहलाता है:
    \begin{choices}\item पक्षपात (bias) \item सटीकता \item ढलान \item औसत\end{choices}
\end{mcq}

\begin{mukhyashabd}
\textbf{अच्छा data} (शुद्धता, विविधता, पर्याप्त मात्रा),
\textbf{पाँच चरण} (संग्रह, सफ़ाई, विश्लेषण, मॉडलन, निष्कर्ष),
\textbf{पक्षपात (bias)}, \textbf{अभिलक्षण (feature)}, \textbf{सहसम्बन्ध (correlation)}।
\end{mukhyashabd}
\selfcheckqr{https://kkreboot.github.io/ai-in-science-hindi/quiz/ch5.html}
